\chapter{Hemoglobin}

\section*{What Is This Test?}
Hemoglobin (Hgb or Hb) is an iron-containing protein within red blood cells that transports oxygen from the lungs to peripheral tissues and carries carbon dioxide from tissues back to the lungs for exhalation. The hemoglobin test measures the concentration of hemoglobin in whole blood and is the most direct indicator of the blood's oxygen-carrying capacity. Each hemoglobin molecule is a tetramer consisting of two alpha and two beta globin chains (in adult hemoglobin, HbA), each containing a heme group with an iron atom at its center capable of binding one oxygen molecule. The hemoglobin concentration is the single most important parameter for diagnosing anemia and assessing its severity, as well as for detecting polycythemia. It is used to screen for blood disorders, monitor chronic conditions, evaluate blood loss, and guide transfusion decisions. The World Health Organization defines anemia based on hemoglobin thresholds: below 13.0 g/dL in men and below 12.0 g/dL in non-pregnant women \cite{who2011haemoglobin}.

\section*{Alternative Names}
Hgb; Hb; Hemoglobin concentration; Blood Hemoglobin; Total Hemoglobin; Hb Level

\section*{Principle}
Hemoglobin is measured spectrophotometrically after conversion to a stable derivative \cite{tietz2023textbook}. In the reference cyanmethemoglobin (HiCN) method, whole blood is mixed with Drabkin's reagent containing potassium ferricyanide and potassium cyanide. Ferricyanide oxidizes hemoglobin iron from the ferrous (Fe\textsuperscript{2+}) to the ferric (Fe\textsuperscript{3+}) state to form methemoglobin, which then combines with cyanide to form cyanmethemoglobin. The absorbance of cyanmethemoglobin is measured at 540 nm spectrophotometrically, and hemoglobin concentration is calculated using Beer's law based on the millimolar absorptivity of cyanmethemoglobin (11.0 L/mmol/cm). In modern automated hematology analyzers (Sysmex XN series, Beckman Coulter DxH, Siemens ADVIA), hemoglobin is measured using a non-cyanide method employing sodium lauryl sulfate (SLS) as the lysing agent, which converts hemoglobin to SLS-methemoglobin measured at 540--555 nm \cite{bain2017dacie}. This eliminates the hazardous cyanide waste issue. Some point-of-care devices (HemoCue) use the azide-methemoglobin method with absorbance measurement at 565 nm and 880 nm for turbidity correction.

\section*{History}
The oxygen-carrying property of blood was demonstrated in the 17th century, but hemoglobin itself was not identified until Felix Hoppe-Seyler crystallized it in 1862 and named it ``h\"{a}moglobin'' \cite{schechter2008hemoglobin}. The first quantitative method for hemoglobin was developed by Gowers in 1878 using a visual color comparator (the Gowers hemoglobinometer). The Tallqvist scale, a paper strip color comparison method, was introduced in 1900 and widely used for decades despite its inaccuracy. In the 1920s, Evelyn and Malloy developed the first photoelectric method for hemoglobin. The cyanmethemoglobin method was introduced by Drabkin and Austin in 1932 and established as the international reference method by the International Committee for Standardization in Haematology (ICSH) in 1966 \cite{kaushansky2021williams}. The SLS method was developed in the 1980s as a cyanide-free alternative and is now the dominant method in automated hematology analyzers. Current research focuses on non-invasive hemoglobin measurement using pulse oximetry and spectrophotometric sensors.

\section*{Why This Test Is Ordered}
\begin{itemize}
  \item Screening for anemia --- the most common hematologic disorder worldwide, particularly in women of childbearing age, children, and the elderly
  \item Classification of anemia severity: mild (10 g/dL to lower limit of normal), moderate (8--10 g/dL), severe (6.5--8 g/dL), and life-threatening (<6.5 g/dL)
  \item Diagnosing polycythemia --- hemoglobin above 16.5 g/dL in women or 18.5 g/dL in men requires investigation for polycythemia vera or secondary causes
  \item Monitoring response to iron supplementation, vitamin B12 therapy, or erythropoietin administration in chronic kidney disease
  \item Preoperative evaluation and assessment of surgical risk
  \item Monitoring blood loss during surgery, trauma, or gastrointestinal hemorrhage
  \item Screening for hemoglobinopathies (sickle cell disease, thalassemia) in high-risk populations or newborns
  \item Monitoring hematologic effects of chemotherapy, radiation therapy, or antiretroviral therapy for HIV
  \item Evaluating fatigue, pallor, dyspnea, tachycardia, or other symptoms suggestive of anemia
\end{itemize}

\section*{Is This a Primary or Secondary Test}
Hemoglobin measurement is a primary screening test for anemia and polycythemia and is a core component of the complete blood count \cite{tierney2023current}.

\section*{How This Test Is Performed}
A venous blood sample is collected into a lavender-top (EDTA) tube from a peripheral vein, typically in the antecubital fossa. A tourniquet is applied briefly, the skin is cleaned with an antiseptic, and a needle is inserted to draw 2--5 mL of blood. The tube is gently inverted 8--10 times to ensure adequate mixing with EDTA anticoagulant. The sample is aspirated by the automated hematology analyzer, which lyses red blood cells and measures hemoglobin spectrophotometrically. Fingerstick capillary blood can also be used with point-of-care hemoglobinometers (HemoCue, Masimo). The entire venipuncture procedure takes less than 5 minutes.

\section*{What Preparation Is Needed}
No special preparation is required. Fasting is not necessary for hemoglobin measurement alone, though if other tests requiring fasting (glucose, lipid panel) are ordered simultaneously, the patient should fast for 8--12 hours. Patients should avoid vigorous exercise for 30 minutes prior to testing, as intense exercise can transiently elevate hemoglobin due to plasma volume contraction and splenic contraction. The patient should be seated or supine for at least 5 minutes before venipuncture, as postural changes can alter plasma volume and hemoglobin concentration \cite{wallach2022interpretation}.

\section*{Contraindications and Precautions}
No absolute contraindications. Factors affecting results include: dehydration (falsely elevated hemoglobin due to hemoconcentration), pregnancy (physiological decrease of 1--2 g/dL due to plasma volume expansion, reaching a nadir at 30--34 weeks), smoking (carbon monoxide increases hemoglobin by 0.5--1.5 g/dL), high altitude (>1,500 meters causes compensatory increase), prolonged tourniquet application (>1 minute causes hemoconcentration with 5--10\% false elevation), severe lipemia (turbidity may interfere with spectrophotometric measurement), very high white blood cell count (>100,000/$\mu$L may falsely elevate hemoglobin), and recent blood transfusion (alters hemoglobin concentration). Plasma hemoglobin from in-vitro hemolysis (traumatic venipuncture, small-gauge needle, vigorous shaking) invalidates the result. The EDTA sample should be analyzed within 4 hours at room temperature or 24 hours refrigerated at 2--8\textdegree{}C.

\section*{Risks}
The risks of venipuncture are minimal and include: mild pain or stinging at the needle insertion site, bruising or small hematoma at the puncture site, lightheadedness or vasovagal syncope (fainting), and extremely rarely, infection, nerve injury, or excessive bleeding in patients with coagulopathy. Fingerstick sampling carries a slight risk of local infection and inaccurate results if the finger is squeezed excessively (diluting the sample with tissue fluid).

\section*{What Happens After the Test}
Hemoglobin results are typically available within 30--60 minutes. The healthcare provider interprets the hemoglobin value in the context of the full CBC, reticulocyte count, and red cell indices. If anemia is detected, the MCV and RDW are used to classify the type: microcytic (iron deficiency, thalassemia), normocytic (chronic disease, acute blood loss, hemolysis), or macrocytic (B12/folate deficiency, myelodysplasia). If hemoglobin is dangerously low (<7 g/dL in hemodynamically stable patients, or <8 g/dL in patients with cardiovascular disease), transfusion of packed red blood cells may be considered. If polycythemia is detected, JAK2 mutation testing, serum erythropoietin level, and pulmonary function tests may be ordered.

\section*{Understanding Results}

\subsection*{Below Normal}
\begin{itemize}
  \item \textbf{Iron deficiency anemia:} The most common cause of low hemoglobin worldwide. Characterized by microcytosis (low MCV), hypochromia (low MCH, low MCHC), elevated RDW, low serum ferritin (<30 ng/mL), low transferrin saturation (<16\%), and elevated TIBC. Causes include chronic blood loss (menorrhagia, GI bleeding), inadequate dietary intake, malabsorption (celiac disease, post-gastrectomy), and increased requirements (pregnancy, growth).
  \item \textbf{Anemia of chronic disease:} Normocytic or mildly microcytic anemia associated with chronic inflammation, infection, malignancy, or autoimmune disease. Characterized by low serum iron, low TIBC, normal or elevated ferritin, and elevated hepcidin levels.
  \item \textbf{Thalassemia:} Microcytic hypochromic anemia with normal or elevated RDW. Alpha-thalassemia (deletion of alpha-globin genes) and beta-thalassemia (mutations in beta-globin gene) cause reduced hemoglobin synthesis. Confirmed by hemoglobin electrophoresis showing elevated HbA2 (>3.5\%) in beta-thalassemia trait.
  \item \textbf{Vitamin B12 deficiency:} Macrocytic anemia (MCV >100 fL) with low hemoglobin. Causes include pernicious anemia (autoimmune destruction of parietal cells), gastric bypass surgery, chronic atrophic gastritis, ileal resection, and strict vegan diet. Associated with neurological symptoms (paresthesias, ataxia, memory loss).
  \item \textbf{Folate deficiency:} Macrocytic anemia from inadequate dietary intake (alcoholism, elderly with poor nutrition), malabsorption, increased requirements (pregnancy, hemolytic anemia), or medications (methotrexate, trimethoprim, phenytoin).
  \item \textbf{Hemolytic anemia:} Autoimmune hemolytic anemia (positive direct Coombs test), hereditary spherocytosis (increased MCHC, spherocytes on smear), G6PD deficiency (Heinz bodies, bite cells), sickle cell disease (HbS on electrophoresis, sickled cells), drug-induced hemolysis.
  \item \textbf{Aplastic anemia:} Pancytopenia with markedly low hemoglobin from bone marrow failure, caused by radiation, chemotherapy, benzene, viral hepatitis, or idiopathic.
  \item \textbf{Chronic kidney disease:} Normocytic anemia from insufficient erythropoietin production, typically when GFR <30 mL/min/1.73m\textsuperscript{2}.
  \item \textbf{Acute blood loss:} Trauma, GI bleeding (peptic ulcer, esophageal varices, diverticulosis), ruptured ectopic pregnancy, or postpartum hemorrhage.
  \item \textbf{Myelodysplastic syndromes:} Macrocytic or normocytic anemia with cytopenias and dysplastic bone marrow morphology, primarily in older adults.
\end{itemize}

\subsection*{Above Normal}
\begin{itemize}
  \item \textbf{Polycythemia vera:} Primary myeloproliferative neoplasm with JAK2 V617F mutation causing unregulated erythroid proliferation. Hemoglobin >16.5 g/dL (women) or >18.5 g/dL (men), with elevated RBC mass, low serum erythropoietin, and often thrombocytosis and leukocytosis.
  \item \textbf{Secondary polycythemia:} Chronic hypoxia from COPD, obstructive sleep apnea, cyanotic congenital heart disease, or high-altitude residence. Elevated erythropoietin with appropriate compensatory erythrocytosis.
  \item \textbf{Ectopic erythropoietin production:} Renal cell carcinoma, hepatocellular carcinoma, cerebellar hemangioblastoma, uterine leiomyoma.
  \item \textbf{Dehydration and hemoconcentration:} Severe vomiting, diarrhea, diuretic overuse, burns, or diabetic ketoacidosis causing relative polycythemia.
  \item \textbf{Exogenous androgen use:} Testosterone replacement therapy, anabolic steroid abuse.
  \item \textbf{Carbon monoxide poisoning:} Chronic exposure from smoking or occupational sources elevates carboxyhemoglobin, reducing functional oxygen-carrying capacity and stimulating compensatory erythrocytosis.
  \item \textbf{High-affinity hemoglobin variants:} Rare hemoglobin mutations (Hb Chesapeake, Hb Yakima) with increased oxygen affinity causing tissue hypoxia and compensatory erythrocytosis.
\end{itemize}

\section*{Normal Results}
\begin{itemize}
  \item \textbf{Adult males:} 13.5--17.5 g/dL (135--175 g/L) \cite{fischbach2023manual}
  \item \textbf{Adult females (non-pregnant):} 12.0--16.0 g/dL (120--160 g/L)
  \item \textbf{Pregnant females:} 11.0--14.0 g/dL (first trimester); 10.5--14.0 g/dL (second trimester); 10.0--14.0 g/dL (third trimester)
  \item \textbf{Children (6--12 years):} 11.5--15.5 g/dL
  \item \textbf{Children (1--6 years):} 11.0--14.0 g/dL
  \item \textbf{Infants (1--12 months):} 10.0--14.0 g/dL
  \item \textbf{Newborns (first week):} 14.0--24.0 g/dL (decreases to 10.0--14.0 g/dL by 2--3 months: physiological nadir)
  \item \textbf{WHO anemia criteria:} <13.0 g/dL (men), <12.0 g/dL (non-pregnant women), <11.0 g/dL (pregnant women)
  \item \textbf{High altitude (>2,400 m):} Males 15.0--20.0 g/dL; Females 14.0--18.5 g/dL
\end{itemize}

\section*{Follow-up Tests}
\begin{itemize}
  \item \textbf{Complete blood count with differential and red cell indices:} Provides MCV, MCH, MCHC, RDW, WBC, and platelet count for comprehensive anemia classification
  \item \textbf{Peripheral blood smear:} Microscopic examination for red blood cell morphology --- hypochromia, microcytosis, macrocytosis, target cells, spherocytes, schistocytes, sickle cells, basophilic stippling, or polychromasia
  \item \textbf{Reticulocyte count with reticulocyte production index:} Differentiates hypoproliferative anemia (low reticulocytes) from hemolytic anemia or recovery from blood loss (elevated reticulocytes)
  \item \textbf{Iron studies:} Serum iron, ferritin, total iron-binding capacity (TIBC), and transferrin saturation
  \item \textbf{Vitamin B12 and serum folate levels:} If macrocytic anemia is present (MCV >100 fL)
  \item \textbf{Hemoglobin electrophoresis:} Identifies and quantifies hemoglobin variants (HbA, HbA2, HbF, HbS, HbC) for hemoglobinopathy diagnosis
  \item \textbf{Direct antiglobulin test (Coombs test):} If autoimmune hemolytic anemia is suspected
  \item \textbf{Haptoglobin and lactate dehydrogenase (LDH):} Markers of hemolysis --- low haptoglobin and elevated LDH indicate intravascular hemolysis
  \item \textbf{Serum erythropoietin level:} Distinguishes primary (low EPO) from secondary (high EPO) polycythemia
  \item \textbf{JAK2 V617F mutation analysis:} Diagnostic for polycythemia vera when hemoglobin is persistently elevated
  \item \textbf{Bone marrow aspiration and biopsy:} When aplastic anemia, myelodysplastic syndrome, or leukemia is suspected
  \item \textbf{G6PD enzyme assay:} If G6PD deficiency hemolysis is suspected (bite cells, Heinz bodies on smear)
\end{itemize}

\section*{References}
\bibliographystyle{unsrtnat}
\bibliography{references}

\clearpage
